\chapter{Introduzione}\label{chap:first}

\inlineminitoc

La pubblica amministrazione locale italiana fa quotidianamente affidamento su una vasta e complessa mole di testi normativi, atti e delibere. L'elaborazione automatica di tali documenti è tuttavia ostacolata da diverse sfide intrinseche: i testi sono spesso archiviati in formati non strutturati (come i PDF), presentano difformità di impaginazione, impiegano un linguaggio burocratico peculiare e si articolano in una rigida e profonda struttura gerarchica. 
Quando si tenta di automatizzare la ricerca all'interno di tali documenti mediante sistemi basati su Intelligenza Artificiale Generativa, un approccio fondato esclusivamente sul recupero di frammenti testuali tramite similarità semantica può non rappresentare esplicitamente la struttura gerarchica del documento e le relazioni esistenti tra le diverse disposizioni. Nel dominio giuridico, dove la corretta individuazione del contesto costituisce una componente centrale del processo di retrieval, tale aspetto assume particolare rilevanza \cite{pipitone2024legalbenchrag}.
È in questo scenario che si inserisce il lavoro di tesi, svolto nell'ambito delle attività di tirocinio presso \textbf{InfoCube} per contribuire allo sviluppo del progetto \textbf{LUDOVICA}, il cui obiettivo finale è fornire ai funzionari della PA un motore generativo affidabile, capace di recuperare gli atti pregressi pertinenti e generare template documentali pronti all'uso. Affinché le risposte del sistema siano rigorose e verificabili è emersa l'esigenza di adottare un'architettura dati che introduca una solida topologia informativa.

Il problema tecnico fondamentale affrontato consiste quindi nella necessità di \textbf{convertire i documenti normativi complessi in una base di conoscenza interrogabile} (\textit{Knowledge Graph}) che conservi fedelmente la struttura semantica e gerarchica originale. In letteratura, i Knowledge Graph sono impiegati per rappresentare entità e relazioni in forma esplicita, rendendo interrogabili collezioni di conoscenza eterogenee e interconnesse \cite{hogan2021knowledgegraphs}. Nel dominio giuridico, inoltre, standard come Akoma Ntoso formalizzano la natura gerarchica dei documenti legislativi e amministrativi e prevedono meccanismi espliciti per citazioni e riferimenti incrociati \cite{akoma_ntoso_2018}.

\section{Obiettivi del lavoro di tesi}
Al fine di superare i limiti esposti e dotare il sistema di solidi meccanismi per l'interpretazione del linguaggio giuridico, il lavoro si è articolato nei seguenti obiettivi:
\begin{enumerate}
    \item \textbf{Modellazione topologica e digitalizzazione:} Costruire una pipeline di inserimento per tradurre documenti legali non strutturati in un database a grafo (\textit{Neo4j}), definendo rigorosamente la mappatura delle relazioni.
    \item \textbf{Semantic Search per il dominio giuridico:} Sviluppare il modulo di interrogazione intelligente del database e confrontare il comportamento di più modelli di embedding nel recupero di disposizioni normative in lingua italiana.
    \item \textbf{Analisi degli approcci di Retrieval:} Misurare, mediante un benchmark diagnostico, in quali condizioni una strategia che impiega le relazioni giuridiche del grafo migliori o peggiori il recupero rispetto a una baseline puramente vettoriale basata sulla \textit{cosine similarity}.
\end{enumerate}

Per raggiungere gli obiettivi prefissati, è stata ideata un'architettura di elaborazione in grado di trasformare documenti non strutturati in una rete di nodi interconnessi. L'approccio proposto si articola in una pipeline a più fasi: inizialmente, i documenti originali (principalmente in formato PDF) vengono processati tramite strumenti di parsing avanzati e OCR per estrarne accuratamente il testo e convertirlo in un formato intermedio (Markdown). Successivamente, un \textit{Large Language Model} analizza tale testo per riconoscerne la suddivisione gerarchica (Titoli, Capi, Articoli, Commi) e le relazioni intrinseche tipiche della struttura giuridica (ad esempio un comma che cita un altro comma).
Questi elementi vengono strutturati in formato JSON e iniettati all'interno del grafo, in cui ogni entità diventa un nodo e l'appartenenza a una data sezione diventa una relazione. Successivamente, i contenuti testuali vengono vettorializzati e salvati sotto forma di \textit{embedding}.
In fase di interrogazione, il sistema affianca alla similarità semantica la navigazione esplicita delle relazioni del grafo, con l'obiettivo di promuovere disposizioni collegate che completano il contesto recuperato. Il presente lavoro valuta il recupero dei commi e il loro ordinamento; la qualità di un'eventuale risposta generata a valle rimane fuori dal perimetro sperimentale.

\section{Confronto con tecniche esistenti}
Le architetture di \textit{Retrieval-Augmented Generation} (RAG) combinano le capacità generative dei \textit{Large Language Models} con il recupero di informazioni provenienti da una base di conoscenza esterna \cite{lewis2020rag}. Nelle implementazioni basate su retrieval vettoriale, i documenti vengono generalmente suddivisi in porzioni testuali, rappresentate mediante embedding e recuperate sulla base della similarità semantica rispetto alla query dell'utente.

L'efficacia di questo processo dipende tuttavia non soltanto dal modello generativo, ma anche dalla qualità delle informazioni recuperate, tale aspetto assume particolare rilevanza nel dominio giuridico. Pipitone e Houir Alami. \cite{pipitone2024legalbenchrag}, attraverso il benchmark \textit{LegalBench-RAG}, evidenziano l'importanza di recuperare segmenti testuali minimi e altamente pertinenti, mostrando come il recupero di porzioni di testo troppo ampie o imprecise possa introdurre informazioni non necessarie nel contesto fornito al modello.

Un ulteriore limite deriva dal fatto che la suddivisione di un documento in frammenti indipendenti tende a non rappresentarne esplicitamente l'organizzazione strutturale. Questo aspetto è particolarmente rilevante per i testi normativi, nei quali una disposizione non costituisce necessariamente un'unità isolata, ma appartiene a una gerarchia composta, ad esempio, da titoli, capi, articoli e commi e può contenere riferimenti ad altre disposizioni. Il solo valore di similarità tra la query e il contenuto testuale consente quindi di stimare la pertinenza semantica di un frammento, ma non codifica direttamente tali relazioni strutturali.

Per affrontare problemi di questo tipo sono stati proposti approcci che affiancano al retrieval vettoriale rappresentazioni strutturate della conoscenza. Edge et al. \cite{edge2024graphrag} propongono un approccio \textit{GraphRAG} nel quale, a partire dai documenti del corpus, viene costruito un Knowledge Graph di entità e relazioni. Il grafo viene successivamente organizzato in entità e relazioni, in questo modo, il sistema può recuperare informazioni distribuite in differenti parti del corpus e affrontare interrogazioni che richiedono una visione più globale dei documenti.

Alla luce di tali considerazioni, il sistema sviluppato in questa tesi conserva esplicitamente la struttura nativa degli atti normativi e alcune delle relazioni presenti tra i relativi elementi, affiancandole alla rappresentazione vettoriale del contenuto testuale. L'ipotesi alla base del lavoro è che tali informazioni strutturali possano contribuire al recupero del contesto normativo rilevante. La superiorità dell'approccio basato sul grafo non viene pertanto assunta a priori, ma viene valutata sperimentalmente attraverso il confronto con una baseline esclusivamente vettoriale.

\section{Struttura della tesi}
Il presente elaborato è organizzato in sei capitoli, di seguito brevemente introdotti.

Il \textbf{Capitolo 2 — Tecnologie utilizzate} fornisce una panoramica dello stack tecnologico adottato: dai linguaggi di programmazione agli strumenti di parsing e OCR per l'estrazione del testo, passando per il database a grafo, fino ai modelli di linguaggio e di embedding impiegati per la strutturazione semantica e la vettorializzazione dei contenuti normativi.

Il \textbf{Capitolo 3 — Architettura del Sistema} illustra l'architettura complessiva del sistema, descrivendo i componenti principali e il flusso di dati attraverso le diverse fasi della pipeline: dall'acquisizione dei documenti grezzi fino all'interrogazione da parte degli utenti finali.

Il \textbf{Capitolo 4 — Implementazione Software} approfondisce nel dettaglio le singole fasi della pipeline di elaborazione: l'estrazione del testo dai documenti PDF, la strutturazione gerarchica operata dai Large Language Models, l'ingestion dei dati strutturati su Neo4j e la creazione degli embedding per la ricerca semantica.

Il \textbf{Capitolo 5 — Ricerca Semantica} definisce le due strategie di recupero confrontate nel lavoro: la baseline puramente vettoriale e la strategia ibrida che usa le relazioni giuridiche come segnale di riordinamento. Ne descrive inoltre i componenti, i parametri e le scelte progettuali.

Il \textbf{Capitolo 6 — Valutazione Sperimentale} presenta il benchmark diagnostico, le metriche e i risultati del confronto, distinguendo i casi in cui la struttura migliora la completezza del recupero da quelli in cui introduce un costo di ordinamento o di selezione.
